\section{Threat Model and Assumptions}
\label{sec:threat}

This section defines the threat model, the role of the LLM attacker, the defender assumptions, and the safety boundaries enforced by CPSForge.

\subsection{Adversary Goals}

The adversary aims to manipulate the cyber-physical process in a way that causes measurable impact while remaining process-valid and, when possible, difficult to detect. Depending on the scene, impact may include unsafe state transitions, degraded performance, incorrect sequencing, product handling failures, or violations of expected operational behavior. The adversary may also seek to delay or evade detection by maintaining plausible sensor and actuator patterns.

\subsection{Adversary Capabilities}

The attacker in CPSForge models a post-compromise adversary who has achieved S7 protocol-level network access to the PLC control network --- for example, through credential theft, insider access, exploitation of an engineering workstation, or lateral movement from the IT network. This access level is realistic: S7 communication between PLC and SCADA/HMI systems is commonly exposed on flat industrial networks, and tools such as python-snap7 allow any network-reachable host to read and write PLC data blocks without authentication on many deployed systems.

With S7-level access, the adversary can: (i)~continuously read the live tag stream (sensor values, actuator states, setpoints) at polling rate; (ii)~infer the current operational phase from tag value patterns; (iii)~inject writes to any writable tag in the PLC data block. Rather than modeling raw memory addresses or firmware-level access, CPSForge expresses attacks at the process-semantic level --- the same level at which a process engineer or SCADA operator would interact with the system. The exact attack surface (which tags are writable by the attacker) is scene-specific and explicitly configured in the scene YAML.

Within this scope, the attacker can perform actuator overrides, setpoint shifts, sensor spoofing, timing delays, and sequence perturbations --- all expressed as structured writes to the PLC tag stream.

\subsection{LLM Role}

The LLM serves as a man-in-the-middle agent in the PLC communication stream, not as a trusted controller. Every three polling cycles (1.5\,s), the LLM receives a structured snapshot of the observed tag stream (whose richness depends on the context tier) and returns a structured decision. As red team: \emph{attack} (with a complete write specification: tag, value, duration, action type) or \emph{wait}. As blue team: \emph{allow} or \emph{block} a proposed write, with a suspicion score and reasoning.

The LLM does not directly access PLC memory, bypass process interlocks, or issue unrestricted writes. All LLM outputs are treated as untrusted proposals until validated by the safety shield. The shield enforces tag whitelists, value ranges, duration caps, and process invariants before any write reaches the PLC.

This architecture separates LLM reasoning capability (which we evaluate) from execution privilege (which is strictly bounded). The paper evaluates how useful LLM communication-analysis outputs are under varying context depths (RQ1), whether training on real communication traces improves those outputs (RQ2), and whether online stream feedback is necessary (RQ3).

\subsection{Timing Model and Latency Characterization}
\label{sec:threat:timing}

CPSForge does \emph{not} operate at the PLC scan-cycle timescale (${\sim}$1--10\,ms). LLM inference on the 3B model takes 13--33\,s per call depending on prompt length and context tier. Each LLM call is issued every three polling cycles (1.5\,s of observation time); during inference, polling continues but no new attacker decisions are made until the call returns. Total wall-clock time per 60-step run is 6--10~minutes.

This latency profile positions the LLM adversary as a \emph{supervisory-timescale} attacker---analogous to an operator or SCADA engineer making strategic process interventions at multi-second intervals---not a sub-millisecond PLC-level manipulator. This is a realistic and well-documented threat model: documented ICS incidents operate on timescales far longer than the PLC scan cycle. The Stuxnet campaign manipulated centrifuge speeds over weeks~\cite{langner2011stuxnet}, the 2015 Ukrainian grid attack coordinated breaker openings over approximately 30~minutes~\cite{liang2017ukraine}, and the TRITON/TRISIS attack on safety instrumented systems required multi-hour reconnaissance before a single write~\cite{johnson2017triton}. Even the fastest known ICS MitM attacks (e.g., HARVEY~\cite{garcia2017harvey}) operate at control-loop timescales (${\sim}$100\,ms--1\,s), not PLC scan-cycle speeds.

Our 30-second observation window with up to 10 strategic writes is conservative relative to real-world ICS attack timelines, where adversaries typically have persistent access over hours or days. The 10-attack budget tests whether the LLM can \emph{select the right writes} from limited opportunities---a harder task than unrestricted injection. This budget also keeps experiments within Factory~I/O's process cycle time (most sorting sequences complete in 20--40\,s), ensuring each run observes at least one complete operational cycle.

\subsection{Defender Model}

The defender has access to the observed plant state, controller state, alarms, and experiment telemetry collected from the PLC-process loop. Depending on the detector configuration, the defender may use static thresholds, process invariants, or learned sequence models. The defender does not assume oracle knowledge of the attacker’s true intent, though logged attacks can later be used for offline analysis and adaptation.

The defender is evaluated on its ability to detect process deviations, unsafe states, suspicious sequences, or other operational anomalies caused by the attacker.

\subsection{Safety Boundary and Trusted Components}

CPSForge assumes that the safety shield is part of the trusted computing base for experimentation. The shield is responsible for validating every candidate action before it reaches the PLC. In practice, it enforces writable-tag whitelists, legal value ranges, maximum override durations, cooldown periods, forbidden simultaneous actions, and scene-specific safety invariants.

The real PLC, deployed logic, and scene configuration are also assumed to be correctly installed and maintained by the operator. Human supervision remains in scope, and the framework assumes that emergency-stop and reset procedures are available during live experiments.

\subsection{Scope and Limitations}

The present work does not model arbitrary firmware compromise, PLC root access, or unrestricted malicious code deployment. The attacker operates within a configured action space mediated by the safety shield, which bounds all writes to process-safe ranges. This bounded model is a deliberate design choice: it enables safe, repeatable, and auditable experimentation while still permitting the attacker to cause measurable process impact---as the evaluation demonstrates. However, this means experimental attack success rates reflect effectiveness \emph{within the configured envelope}, not the full capability of an unconstrained adversary. The framework is designed to study whether LLMs can produce process-grounded attack strategies and whether defenders can improve under such pressure; it does not claim to replicate the full scope of nation-state ICS attack campaigns.

\subsection{Extended Threat Model: Three Attack Tiers}
\label{sec:threat:tiers}

While the original evaluation focuses on I/O-level data manipulation (Tier~1), we extend the threat model to capture the full spectrum of PLC-level attacks observed in real-world incidents. This multi-tier model positions CPSForge's contributions within a broader taxonomy and motivates two additional testbeds.

\subsubsection{Tier 1: I/O Data MitM (Original Evaluation)}
The attacker reads and writes PLC I/O tags (sensors, actuators, setpoints) via S7 protocol. Written values are \emph{overwritten by PLC logic each scan cycle}, requiring continuous or precisely timed injection. This tier maps to the Online MITM Attacker (\S\ref{sec:design}) and covers all 333~runs in the evaluation. Real-world analogues include Oldsmar (single setpoint write), BlackEnergy (breaker commands via SCADA HMI), and Industroyer (automated protocol-level commands)---see Table~\ref{tab:incidents} in \S\ref{sec:background:incidents}.

\subsubsection{Tier 2: Deep State Manipulation (New)}
The attacker targets \emph{internal PLC state variables}---state machine indices (\texttt{iState}), timers, classification flags, thresholds, and enable bits---that are accessible via the same S7 protocol but are \textbf{never refreshed by PLC logic}. Unlike Tier~1 writes that the PLC immediately overwrites, a single Tier~2 write to a variable like \texttt{Enable} (BOOL) or \texttt{LightThresh} (REAL) persists indefinitely until the scene resets, because the PLC reads but never re-initializes these variables during normal operation.

This tier is analogous to Stuxnet's writes to DB890, which modified PLC behavior by changing data block values rather than modifying the PLC program itself~\cite{langner2013kill}. The key insight is that \emph{behavioral modification} of a PLC can be achieved through data alone---by targeting internal state variables that the control logic trusts implicitly.

We implement this tier as the \emph{DeepStateMITMAttacker}, which shares the same observation--phase--context--decide--act loop as the Tier~1 attacker but operates on an extended attack surface (\texttt{deep\_state\_attack\_surface} in the scene YAML) with eight attack categories: threshold manipulation, state machine jump, timer acceleration, safety kill switch, decision flag spoofing, setpoint hijack, counter spoofing, and mode manipulation.

The primary defense for this tier is the \emph{StateConsistencyChecker}, which enforces state transition validity (e.g., \texttt{iState} can only advance to adjacent states), timer progression bounds, protected variable guards (e.g., \texttt{Enable} is never writable during operation), and classification flag phase-locking.

\subsubsection{Tier 3: Logic Code Generation (New)}
The attacker analyzes PLC control logic (SCL source code) and generates adversarial modifications---minimal, compilable code changes designed to disrupt the process while remaining difficult to detect through code review. This tier requires engineering workstation access for deployment (manual TIA Portal upload) and tests the LLM's \emph{cognitive capability} to comprehend and weaponize structured control logic.

This maps directly to Stuxnet's offline payload development methodology: the attack payload (modified OB1/OB35) was crafted by analyzing centrifuge control logic, then deployed via USB~\cite{falliere2011w32}. The LLM's contribution is the same \emph{analysis and generation} step; deployment is out of scope.

The \emph{LogicAnalyzer} scans SCL source for seven vulnerability categories (dead code injection, comparison inversion, timer alteration, interlock removal, gradual drift, state shortcut, sensor clamping) and generates minimal adversarial code modifications with stealth assessments and detection difficulty ratings. The symmetric defense, \emph{CodeReviewDefender}, reviews SCL diffs using the same LLM architecture and outputs approve/reject decisions with suspicion scores.

\subsubsection{LLM-Accessible Attack Surface Enumeration}
\label{sec:threat:attack_surface}

A central question for LLM-enabled CPS attacks is: \emph{what can the model actually reach?} Table~\ref{tab:attack_surface} enumerates the concrete variables accessible via S7 protocol across our three evaluation scenes. All variables reside in non-optimized PLC data blocks (DBs), which python-snap7 can read and write by byte offset without authentication on default S7-1200 configurations.

\begin{table}[h]
\centering
\caption{LLM-accessible attack surface per scene. ``PLC refresh'' indicates whether the PLC overwrites the variable each scan cycle ($\checkmark$) or trusts the stored value ($\times$).}
\label{tab:attack_surface}
\small
\begin{tabular}{@{}llllc@{}}
\toprule
\textbf{Scene} & \textbf{Variable} & \textbf{Type} & \textbf{Attack Category} & \textbf{PLC Refresh} \\
\midrule
\multirow{4}{*}{\shortstack[l]{Level\\Control}} & SetpointIn & REAL & Setpoint hijack & $\times$ \\
& FillValve & REAL & Actuator override & $\checkmark$ \\
& DischargeValve & REAL & Actuator override & $\checkmark$ \\
& Enable & BOOL & Safety kill & $\times$ \\
\midrule
\multirow{5}{*}{\shortstack[l]{Sorting\\Weight}} & LightThresh & REAL & Threshold manip. & $\times$ \\
& HeavyThresh & REAL & Threshold manip. & $\times$ \\
& iState & INT & State machine jump & $\checkmark$ \\
& Conveyor1/2/3 & BOOL & Actuator override & $\checkmark$ \\
& Sorter1/2 & BOOL & Actuator override & $\checkmark$ \\
\midrule
\multirow{5}{*}{\shortstack[l]{Sorting\\Height}} & iState & INT & State machine jump & $\checkmark$ \\
& TurnTimer & INT & Timer acceleration & $\checkmark$ \\
& DwellTimer & INT & Timer acceleration & $\checkmark$ \\
& Enable & BOOL & Safety kill & $\times$ \\
& IsTall & BOOL & Decision flag spoof & $\checkmark$ \\
\bottomrule
\end{tabular}
\end{table}

The distinction between PLC-refreshed ($\checkmark$) and non-refreshed ($\times$) variables is critical. \textbf{Tier~1 attacks} target refreshed variables, requiring the attacker to re-inject each scan cycle or time the injection precisely between PLC writes. \textbf{Tier~2 attacks} target non-refreshed variables, where a single write persists until explicit reset. This analysis demonstrates that even simple PLC programs contain \emph{semantic trust assumptions}---variables like \texttt{SetpointIn}, \texttt{Enable}, and classification thresholds are read by the control logic on every scan but are never re-initialized, creating persistent write targets. This finding generalizes to real-world PLCs, where data block parameters (alarm limits, recipe values, PID tuning constants) are typically writable via the same protocol used for process data, with no access control distinction~\cite{mclaughlin2016cps_integrity}.

\subsubsection{Tier Comparison}

\begin{table}[h]
\centering
\caption{Three-tier threat model comparison}
\label{tab:threat-tiers}
\small
\begin{tabular}{@{}lccc@{}}
\toprule
\textbf{Property} & \textbf{Tier 1: I/O MitM} & \textbf{Tier 2: Deep State} & \textbf{Tier 3: Logic Gen} \\
\midrule
Protocol & S7 (python-snap7) & S7 (python-snap7) & TIA Portal (manual) \\
Target & I/O tags & Internal state & PLC program \\
Persistence & 1 scan cycle & Permanent & Permanent \\
PLC overwrites? & Yes & No & N/A (is the logic) \\
Automation & Fully automated & Fully automated & Manual upload \\
Real-world & Oldsmar, Industroyer & Stuxnet DB890 & Stuxnet OB1/OB35 \\
Defense & Safety shield + LLM & State consistency & Code review \\
\bottomrule
\end{tabular}
\end{table}